% =============================================================================
% TCC II — Monografia
% UFOPA / IEG — Bacharelado em Ciência da Computação
%
% Mesmo preâmbulo do TCC I (../comum/ufopa-tcc.sty) e mesma bibliografia
% (../comum/refs.bib). O que muda: natureza do trabalho, ano, elementos
% pré-textuais adicionais e as seções.
% =============================================================================

\documentclass[12pt,a4paper,oneside]{article}

\usepackage{ufopa-tcc}

% --- Dados do trabalho -------------------------------------------------------
\ufopatitulo{Sistemas inteligentes aplicados à gestão pública}
\ufopasubtitulo{automação da fase preparatória e pesquisa de preços em
  licitações municipais}
\ufopaautor{Rafael Nobre de Souza}
\ufopaorientador{Prof.\textsuperscript{a} \pendente{titulação: Dra.\ ou Ma.?}
  Socorro Vânia Lourenço Alves}
\ufopainstituicao{Universidade Federal do Oeste do Pará}
\ufopaunidade{Instituto de Engenharia e Geociências}
\ufopacurso{Bacharelado em Ciência da Computação}
\ufopalocal{Santarém --- Pará}
\ufopaano{2027}
\ufopanatureza{Trabalho de Conclusão de Curso apresentado como requisito para
  obtenção do grau de Bacharel em Ciência da Computação, pela Universidade
  Federal do Oeste do Pará, no Instituto de Engenharia e Geociências.}

\begin{document}

% --- Elementos pré-textuais ---------------------------------------------------
\ufopainiciopretextual
\ufopacapa
\ufopafolhaderosto
\ufopafichacatalografica{figuras/ficha-catalografica.pdf}
\ufopafolhadeaprovacao{%
  \pendente{preencher após a defesa}%
}{%
  \ufopamembrobanca{\ufopaimprimirorientador}{\ufopaimprimirinstituicao}
  \ufopamembrobanca{\pendente{nome do segundo examinador}}{\pendente{instituição}}
  \ufopamembrobanca{\pendente{nome do terceiro examinador}}{\pendente{instituição}}
}

% Dedicatória, agradecimentos e epígrafe são opcionais (guia 3.2.5 a 3.2.7).
% Descomente e escreva quando fizer sentido — não deixe placeholder na entrega.
% \begin{ufopadedicatoria}...\end{ufopadedicatoria}
% \begin{ufopaagradecimentos}...\end{ufopaagradecimentos}
% \begin{ufopaepigrafe}...\end{ufopaepigrafe}

\ufoparesumo{%
  \pendente{Resumo definitivo, entre \ufoparesumomin{} e \ufoparesumomax{}
  palavras: objetivo, método, resultados obtidos e conclusão.}%
}{%
  \pendente{entre \ufopapalavraschavemin{} e \ufopapalavraschavemax{} termos}%
}
\ufopaabstract{%
  \pendente{Tradução do resumo definitivo.}%
}{%
  \pendente{tradução das palavras-chave}%
}
\ufopasumario

% --- Elementos textuais -------------------------------------------------------
\ufopainiciotextual
\section{Introdução}

\pendente{Evoluir a introdução do TCC I, incorporando o que a execução do
projeto revelou. Ao migrar texto, reconferir cada citação: afirmação que era
"esperada" em TCC I precisa virar afirmação sustentada por resultado em TCC II,
ou ser removida.}

\section{Fundamentação teórica}

\pendente{Consolidar a partir dos fichamentos do vault. Toda afirmação não
trivial precisa de fonte rastreável até um documento lido — não até uma
lembrança de leitura.}

\section{Trabalhos relacionados}

\pendente{Resultado do levantamento sistemático, com o quadro comparativo e a
lacuna que este trabalho ocupa.}

\section{Materiais e métodos}

\pendente{Descrever o que foi efetivamente feito, no passado, e não o que se
pretendia fazer. Divergências em relação ao planejado no TCC I são material
legítimo do texto: registre a mudança e o motivo.}

\subsection{Arquitetura da solução}

\pendente{Diagrama e descrição dos componentes.}

\subsection{Modelagem de dados}

\pendente{Modelo conceitual e lógico, incluindo o mecanismo de trilha de
auditoria que permite reproduzir uma pesquisa de preços na data em que foi
realizada.}

\subsection{Integração com dados públicos de contratações}

\pendente{Endpoints consumidos, parâmetros de consulta, tratamento de falhas e
estratégia de cache.}

\subsection{Tratamento estatístico dos preços}

\pendente{Critério de exclusão de valores discrepantes e cálculo dos
parâmetros, com a justificativa normativa de cada escolha.}

\subsection{Geração assistida de minutas}

\pendente{Como os documentos são produzidos e como sua conformidade é
verificada.}

\subsection{Procedimento de avaliação}

\pendente{Participantes, tarefas, métricas e instrumentos. Detalhar o suficiente
para que outra pessoa possa repetir o experimento.}

\section{Resultados e discussão}

\pendente{Apresentar os dados obtidos antes de interpretá-los. Números sem
medição não entram — se um resultado não foi medido, ele não existe nesta
seção.}

\section{Considerações finais}

\pendente{Retomar o problema e os objetivos, dizendo objetivamente quais foram
atingidos e quais não foram.}

\subsection{Contribuições}

\pendente{O que este trabalho entrega que antes não existia.}

\subsection{Limitações}

\pendente{Escrever com honestidade. Banca reconhece limitação omitida, e a
omissão custa mais caro que a limitação.}

\subsection{Trabalhos futuros}

\pendente{Desdobramentos concretos, não lista de desejos.}


% --- Elementos pós-textuais ---------------------------------------------------
\ufopareferencias

\end{document}
